\subsection{\texorpdfstring{INJECTION OF \(k[\mathbf{P}]\) INTO \(\mathbf{S}(\mathfrak{h})^{*}\)}{INJECTION OF k{[}\textbackslash mathbf\{P\}{]} INTO \textbackslash mathbf\{S\}(\textbackslash mathfrak\{h\})\^{}\{*\}}}

The map \(p \mapsto \exp p\) from P to \(\mathbf{S}(\mathfrak{h})^{*}\) is a homomorphism from the additive group P to \(\mathbf{S}(\mathfrak{h})^{*}\) equipped with its multiplicative structure (no. 1). Consequently, there exists a unique homomorphism \(\psi\) from the algebra k{[}P{]} of the monoid P to the algebra \(\mathbf{S}(\mathfrak{h})^{*}\) such that

\[
\psi (e ^ {\lambda}) = \exp (\lambda) \quad (\lambda \in \mathrm{P})
\]

(in the notations of §7, no. 7). By no. 1, \(\psi\) is injective. By transport of structure, \(\psi(w(e^{\lambda})) = w(\psi(e^{\lambda}))\) for all \(\lambda \in \mathrm{P}\) and all \(w \in \mathrm{W}\) . Hence, if \(k[\mathrm{P}]^{\mathrm{W}}\) (resp. \(\mathbf{S}(\mathfrak{h})^{*\mathrm{W}}\) ) denotes the set of elements of \(k[\mathrm{P}]\) (resp. \(\mathbf{S}(\mathfrak{h})^{*}\) ) invariant under \(\mathrm{W}\) , we have \(\psi(k[\mathrm{P}]^{\mathrm{W}}) \subset \mathbf{S}(\mathfrak{h})^{*\mathrm{W}}\) .

PROPOSITION 1. Let \(S^m(\mathfrak{h}^*)^{\mathrm{W}}\) be the set of elements of \(S^m(\mathfrak{h}^*)\) invariant under W. Then \(\mathrm{pr}_m(\psi(k[\mathrm{P}]^{\mathrm{W}})) = \mathbf{S}^m(\mathfrak{h}^*)^{\mathrm{W}}\) .

It is clear from the preceding that \(\mathrm{pr}_{m}(\psi(k[\mathrm{P}]^{\mathrm{W}})) \subset \mathbf{S}^{m}(\mathfrak{h}^{*})^{\mathrm{W}}\) . Every element of \(\mathbf{S}^{m}(\mathfrak{h}^{*})\) is a k-linear combination of elements of the form

\[
\mathrm{pr} _ {m} (\exp \lambda) = (\mathrm{pr} _ {m} \circ \psi) (e ^ {\lambda})
\]

where \(\lambda \in \mathrm{P}\) (Lemma 1). Hence every element of \(\mathbf{S}^m (\mathfrak{h}^*)^{\mathrm{W}}\) is a linear combination of elements of the form

\[
\sum_ {w \in \mathrm{W}} w ((\mathrm{pr} _ {m} \circ \psi) (e ^ {\lambda})) = (\mathrm{pr} _ {m} \circ \psi) \left(\sum_ {w \in \mathrm{W}} w (e ^ {\lambda})\right),
\]

each of which belongs to \(\mathrm{pr}_m(\psi (k[\mathrm{P}]^{\mathrm{W}}))\) .

PROPOSITION 2. Let E be a finite dimensional g-module. Let U(h) = S(h) be the enveloping algebra of h. If \(u \in U(h)\), then

\[
\operatorname{Tr} (u _ {\mathrm{E}}) = \langle \psi (\operatorname{ch} \mathrm{E}), u \rangle .
\]

It suffices to treat the case in which \(u = h_{1} \ldots h_{m}\) with \(h_{1}, \ldots, h_{m} \in h\) . For all \(\lambda \in P\) , let \(d_{\lambda} = \dim E^{\lambda}\) . Then \(\operatorname{ch} E = \sum_{\lambda} d_{\lambda} e^{\lambda}\) , so \(\psi(\operatorname{ch} E) = \sum_{\lambda} d_{\lambda} \exp(\lambda)\) and hence

\[
\begin{array}{r l} \langle \psi (\mathrm{chE}), u \rangle & = \sum_ {\lambda} d _ {\lambda} \langle \exp \lambda , h _ {1} \dots h _ {m} \rangle \\ & = \sum_ {\lambda} d _ {\lambda} \lambda (h _ {1}) \dots \lambda (h _ {m}) \quad (\text {no. 1}) \\ & = \operatorname{Tr} u _ {\mathrm{E}}. \end{array}
\]

COROLLARY 1. Let \(\mathrm{U}(\mathfrak{g})\) be the enveloping algebra of \(\mathfrak{g}\) . Let the homomorphism \(\zeta : \mathrm{U}(\mathfrak{g})^* \to \mathrm{U}(\mathfrak{h})^* = \mathbf{S}(\mathfrak{h})^*\) be the transpose of the canonical injection \(\mathrm{U}(\mathfrak{h}) \to \mathrm{U}(\mathfrak{g})\) . The following diagram commutes

\[
\begin{array}{c c c} \mathcal {R} (\mathfrak {g}) & \stackrel {{\text {ch}}} {{\longrightarrow}} & \mathbf {Z} [ \mathrm{P} ] \\ \operatorname{Tr} \Big \downarrow & & \Big \downarrow_ {\psi} \\ \mathrm{U} (\mathfrak {g}) ^ {*} & \stackrel {{\zeta}} {{\longrightarrow}} & \mathbf {S} (\mathfrak {h}) ^ {*}. \end{array}
\]

This is simply a reformulation of Prop. 2.

COROLLARY 2. Let m be an integer \(\geq 0\) . Every element of \(\mathbf{S}^{m}(\mathfrak{h}^{*})^{\mathrm{W}}\) is a linear combination of polynomial functions on h of the form \(x \mapsto \operatorname{Tr}(\rho(x)^{m})\) , where \(\rho\) is a finite dimensional linear representation of g.

By Prop. 1, \(\mathbf{S}^m (\mathfrak{h}^*)^{\mathrm{W}} = (\mathrm{pr}_m\circ \psi)(k[\mathrm{P}]^{\mathrm{W}})\) . Now \(\mathbf{Z}[\mathrm{P}]^{\mathrm{W}} = \operatorname {ch}\mathcal{R}(\mathfrak{g})\) (§7, no. 7, Th. 2 (ii)). Thus, by Chap. VI, §3, no. 4, Lemma 3, \(\psi (k[\mathrm{P}]^{\mathrm{W}})\) is the \(k\) -vector subspace of \(\mathbf{S}(\mathfrak{h})^{*}\) generated by \(\psi (\operatorname {ch}\mathcal{R}(\mathfrak{g})) = \zeta (\operatorname {Tr}\mathcal{R}(\mathfrak{g}))\) . Consequently, \(\mathbf{S}^m (\mathfrak{h}^*)^{\mathrm{W}}\) is the vector subspace of \(\mathbf{S}^m (\mathfrak{h}^*)\) generated by \((\mathrm{pr}_m\circ \zeta \circ \mathrm{Tr})(\mathcal{R}(\mathfrak{g}))\) . But, if \(\rho\) is a finite dimensional linear representation of \(\mathfrak{g}\) ,

\[
((\operatorname{pr} _ {m} \circ \zeta \circ \operatorname{Tr}) (\rho)) (x) = \left\langle (\zeta \circ \operatorname{Tr}) (\rho), \frac {x ^ {m}}{m !} \right\rangle = \frac {1}{m !} \operatorname{Tr} (\rho (x) ^ {m})
\]

for all \(x \in h\) .

